\documentclass[a4paper,10.5pt]{article}
\usepackage[utf8]{inputenc}
\usepackage[T1]{fontenc}
\usepackage[french]{babel}
\usepackage[left=1cm,right=1cm,top=.5cm,bottom=0cm]{geometry}
\usepackage{enumitem}
\usepackage{hyperref}
\usepackage{xcolor}
\usepackage{titlesec}
\usepackage{fontawesome5}
\usepackage{lmodern}
\usepackage[normalem]{ulem}

% --- Couleurs Nokia (Bleu corporate) ---
\definecolor{primary}{RGB}{18, 65, 145}
\definecolor{secondary}{RGB}{80, 80, 80}

% --- Configuration des liens ---
\hypersetup{
    colorlinks=true,
    linkcolor=primary,
    filecolor=primary,
    urlcolor=primary,
}

% --- Formatage des titres ---
\titleformat{\section}{\large\bfseries\color{primary}\uppercase}{}{0em}{}[\titlerule]
\titlespacing{\section}{0pt}{8pt}{5pt}

% --- Commandes personnalisées ---
\newcommand{\resumeItem}[1]{\item\small{#1}}
\newcommand{\resumeSubheading}[4]{
  \vspace{-1pt}\item
    \begin{tabular*}{0.97\textwidth}[t]{l@{\extracolsep{\fill}}r}
      \textbf{#1} & \textit{\small #2} \\
      \textit{\small #3} & \textit{\small #4} \\
    \end{tabular*}\vspace{-5pt}
}

\begin{document}

% --- EN-TÊTE ---
\begin{center}
    {\Large \textbf{Loïc Aron MBASSI EWOLO}} \\ \vspace{4pt}
    \textbf{\large Stage : Développement Logiciels Temps Réel 5G/6G} \\
    \textit{\small Élève Ingénieur en 2ème année, double diplôme ENSEA} \\ \vspace{4pt}
    \small
    \faMapMarker\ Cergy, France (Mobile) \quad
    \faPhone\ (+33) 7 59 52 33 98 \quad
    \faEnvelope\ \href{mailto:loicaron.mbassiewolo@ensea.fr}{loicaron.mbassiewolo@ensea.fr} \\ \vspace{2pt}
    \faLinkedin\ \href{https://www.linkedin.com/in/mbassi-loic-3942a9246/}{@mbassi-loic} \quad
    \faGithub\ \href{https://github.com/Nameless0l}{@Nameless0l} \quad
    \faGlobe\ \href{https://mbassiloic.tech}{mbassiloic.tech}
\end{center}

\vspace{-8pt}

% --- PROFIL ---
\noindent\small{Élève-ingénieur en double diplôme à l'\textbf{ENSEA} avec une formation riche et variée en \textbf{Mathématiques, Électronique, Informatique et Programmation}. 
J'aime explorer de nouvelles approches et technologies ainsi que relever des défis open source. 
Ce dont témoignent la note de \textbf{19,25/20} en physique obtenue au Baccalauréat S, la 2\textsuperscript{e} place au hackathon IndabaX \textbf{Africa} et mes divers travaux en programmation \textbf{système}, \textbf{temps réel} et \textbf{embarqué}.}

\vspace{4pt}

% --- FORMATION ---
\section{Formation}
\begin{itemize}[leftmargin=*, label={.}, nosep]
    \item \textbf{Double diplôme : École Nationale Supérieure de l'Électronique et de ses Applications (ENSEA)} \hfill \textit{2025 -- Présent} \\
    \small{2ème année - Spécialité Intelligence Artificielle et \textbf{Traitement du Signal}, Ingénieur de conception}
    \vspace{2pt}
    \item \textbf{École Nationale Supérieure Polytechnique de Yaoundé — Génie Informatique} \hfill \textit{2021 -- Présent} \\
    \small{Niveaux 4 \& 5 : Systèmes d'exploitation, architectures distribuées, programmation concurrente, réseaux} \\
    \small{Niveaux 1 à 3 : Logique, algorithmique, programmation \textbf{C}, électronique | \textbf{Licence 1} Informatique en parallèle}
\end{itemize}

% --- COMPÉTENCES TECHNIQUES ---
\section{Compétences Techniques}
\begin{itemize}[leftmargin=*, nosep]
    \item \small{\textbf{Système \& Temps Réel :} \textbf{Langage C} (IPC, signaux UNIX, mémoire partagée), Linux/Unix, Makefiles, debugging bas-niveau}
    \item \small{\textbf{Embarqué :} Arduino, Raspberry Pi, Nvidia Jetson, TensorFlow Lite, STM32, capteurs IMU/Flex}
    \item \small{\textbf{Outils \& DevOps :} Git/GitHub, \textbf{Docker}, CI/CD, tests unitaires \& intégration, méthodologie \textbf{Agile/Scrum}}
    \item \small{\textbf{IA/ML :} Python, PyTorch, TensorFlow, OpenCV, YOLO | Optimisation modèles pour hardware limité}
    \item \small{\textbf{Développement :} Java, PHP/Laravel, JavaScript, TypeScript | APIs RESTful, Microservices}
\end{itemize}

% --- PROJETS SYSTÈME & EMBARQUÉ ---
\section{Projets — Programmation Système \& Embarqué}
\begin{itemize}[leftmargin=*, label={}]

    \item \textbf{Système de simulation taxi multiprocessus en C — PROJET-TAXI} \hfill \textit{2023 -- 2024, académique}
    \vspace{-2pt}
    \begin{itemize}[leftmargin=0.4cm, nosep, label=\tiny$\bullet$]
        \item \small{Architecture distribuée simulant clients, motos-taxis et centre opérationnel avec contraintes \textbf{temps réel}}
        \item \small{Communication inter-processus : files de messages, \textbf{signaux UNIX}, \textbf{mémoire partagée} (SHM)}
        \item \small{Développement d'un ordonnanceur central gérant l'allocation ressources avec politique FIFO et prévention deadlocks}
        \item \small{Visualisation temps-réel avec \textbf{OpenGL}, Makefiles modulaires, profilage et debugging bas-niveau — \href{https://github.com/Nameless0l/PROJET-TAXI}{\faGithub}}
    \end{itemize}
    \vspace{4pt}

    \item \textbf{Upgrade Jetson — Déploiement IA Embarquée (Challenge CoHoMa, Défense)} \hfill \textit{2025, académique}
    \vspace{-2pt}
    \begin{itemize}[leftmargin=0.4cm, nosep, label=\tiny$\bullet$]
        \item \small{Mise en fonction des Jetsons pour des robots militaires, intégration et optimisation des modèles d'IA légers (YOLOv10)}
        \item \small{Déploiement \textbf{Docker} sur systèmes embarqués, traitement données Lidar 3D VLP16, caméras de profondeur}
        \item \small{Navigation autonome basée sur le \textbf{SLAM}, orchestration des modules développés par les autres équipes}
    \end{itemize}
    \vspace{4pt}

    \item \textbf{Harmony Gloves — Gants Connectés pour Langue des Signes} \hfill \textit{2024, personnel}
    \vspace{-2pt}
    \begin{itemize}[leftmargin=0.4cm, nosep, label=\tiny$\bullet$]
        \item \small{Conception et assemblage de gants instrumentés avec capteurs flex et gyroscopes}
        \item \small{Programmation \textbf{Arduino/STM32} pour lecture données capteurs, fusion avec modèle Computer Vision}
        \item \small{Présenté aux Journées Portes Ouvertes de l'ENSPY 2024}
    \end{itemize}

\end{itemize}

% --- EXPÉRIENCES PROFESSIONNELLES ---
\section{Expériences Professionnelles \& Recherche}
\begin{itemize}[leftmargin=*, label={}]

    \resumeSubheading
      {Stage de Recherche à distance — MILA}{Juin -- Août 2025}
      {Institut Québécois d'Intelligence Artificielle}{Québec, Canada}
      \resumeItem{Grokking (de OpenAI), Généralisation des réseaux de neurones après sur-apprentissage}
      \begin{itemize}[leftmargin=0.4cm, nosep, label=\tiny$\bullet$]
        \item \small{Reproduction d'expériences SOTA, implémentation de pipelines de tests sous \textbf{PyTorch}}
      \end{itemize}
    \vspace{2pt}

    \resumeSubheading
      {Stage de Recherche — Laboratoire IoT, ENSPY}{Oct. 2023 -- Jan. 2024}
      {Supervisé par Dr. Anne Chana}{Yaoundé, Cameroun}
      \resumeItem{Intégration d'algorithmes de Machine Learning avec des dispositifs IoT}
      \begin{itemize}[leftmargin=0.4cm, nosep, label=\tiny$\bullet$]
        \item \small{Traitement et prise de décision en \textbf{temps réel} sur hardware limité (Arduino, Raspberry Pi)}
        \item \small{Optimisation modèles ML pour contraintes mémoire et consommation énergétique — \textbf{TensorFlow Lite}}
      \end{itemize}
    \vspace{2pt}

    \resumeSubheading
      {Développeur Full-Stack (Fintech) — CSC}{Oct. 2024 -- Août 2025}
      {Douala, Cameroun}{}
      \resumeItem{Conception et développement d'une API fintech sécurisée et d'interfaces de gestion}
      \begin{itemize}[leftmargin=0.4cm, nosep, label=\tiny$\bullet$]
        \item \small{Architecture backend (Laravel/Docker, Redis), \textbf{CI/CD}, tests unitaires \& intégration (TDD)}
        \item \small{Méthodologie \textbf{Agile/Scrum}, revue de code, gestion des accès concurrents}
      \end{itemize}

\end{itemize}

% --- AUTRES PROJETS ---
\section{Autres Projets Significatifs}
\begin{itemize}[leftmargin=*, label={}]

    \item \textbf{GoFind — Architecture Microservices} \hfill \textit{2024, académique}
    \vspace{-2pt}
    \begin{itemize}[leftmargin=0.4cm, nosep, label=\tiny$\bullet$]
        \item \small{Plateforme microservices : Laravel, NestJS, MongoDB, communication inter-services avec RabbitMQ, Docker}
    \end{itemize}
    \vspace{2pt}

    \item \textbf{Laravel API Generator — Génération de Code Assistée par LLM} \hfill \textit{2024 -- Présent, Open Source}
    \vspace{-2pt}
    \begin{itemize}[leftmargin=0.4cm, nosep, label=\tiny$\bullet$]
        \item \small{Conception d'un package pour la génération de projets à partir de diagrammes UML, parsing XML, Python}
    \end{itemize}

\end{itemize}

% --- HACKATHONS & LEADERSHIP ---
\section{Hackathons \& Leadership}
\begin{itemize}[leftmargin=*, nosep]
    \item \small{\textbf{Hackathons :} 2\textsuperscript{e} place IndabaX Africa (2025) | 1\textsuperscript{re} place CITS National \& Mountain Hub Regional (2024)}
\end{itemize}

% --- CERTIFICATIONS ---
\section{Certifications \& Langues}
\begin{itemize}[leftmargin=*, nosep]
    \item \small{Supervised Machine Learning (DeepLearning.AI) | Python for Data Science (IBM)}
    \item \small{\textbf{Langues :} Français (Langue maternelle), Anglais (Communication scientifique)}
\end{itemize}

\end{document}
